\section{Summary}
% ===============================================================

We began with the counting statement contained in the Fundamental Theorem of
Algebra and asked for its analogue for meromorphic functions. Making the question
precise required two things: a space on which infinity is an ordinary point,
which is $\PP$; and a device recording all local multiplicities at once, which is
the divisor.

Local analysis (Section~\ref{sec:prelim}) showed that each point contributes a
single integer. Section~\ref{sec:sphere} constructed $\PP$ and identified its
meromorphic functions with $\CC(z)$. Section~\ref{sec:divisors} assembled the
local integers into the homomorphism $\divv : \CC(z)^\times \to \Div(\PP)$, whose
kernel is $\CC^\times$ and whose image is $\Prin(\PP)$; these facts are the two
exact sequences of Section~\ref{sec:exact}. The central result,
Theorem~\ref{thm:main}, is that $\Prin(\PP) = \ker(\deg)$: the balance of zeros
and poles is the only obstruction. From it we obtained $\Cl(\PP) \cong \ZZ$
(Theorem~\ref{thm:cl}) and, by an explicit computation with transition functions
on the two standard charts, $\Pic(\PP) \cong \ZZ$ (Theorem~\ref{thm:pic}).

As explained in Remark~\ref{rem:genus}, all of this is the genus-zero case of a
theory in which the answers are far richer.

% ===============================================================
